\section{Phần B --- Chứng minh lý thuyết}
\label{sec:phan-b}

Mục này chứng minh sáu mệnh đề của Đề tài 9 \cite[tr.~18]{de}. Mỗi mệnh đề nêu
giả thiết, kích thước các đối tượng, các bước suy luận và kết luận. Các chứng minh
không cần đầu vào của Phần~C. Mọi logarit là logarit tự nhiên. Ví dụ số trong mục
này chỉ để minh họa.

Bản độc lập của mục này nằm ở \texttt{docs/theory/part-b.tex} và biên dịch riêng
được.

%--------------------------------------------------------------------
\subsection{B.1 --- Phương sai tích vô hướng và hệ số chia tỷ lệ}
\label{sec:b1}

\textbf{Mệnh đề.} Cho $d_k\ge 1$ và $q,k\in\R^{d_k}$. Giả sử toàn bộ $2d_k$ thành
phần của $q$ và $k$ độc lập lẫn nhau, mỗi thành phần có kỳ vọng $0$ và phương sai
$1$. Khi đó $\Var(q^\top k)=d_k$. Hệ số $1/\sqrt{d_k}$ giữ phương sai của mỗi điểm
chú ý đã chia tỷ lệ bằng $1$ trong mô hình giả thiết này
\cite[§3.2.1, chú thích 4]{attention}.

\subsubsection{Tính phương sai}

Đặt $Z_i=q_i k_i$ với $1\le i\le d_k$. Từ $\Var(q_i)=\E[q_i^2]-(\E[q_i])^2$ và
$\E[q_i]=0$, suy ra $\E[q_i^2]=1$; tương tự $\E[k_i^2]=1$. Vì $q_i$ và $k_i$ độc
lập,
\[
  \E[Z_i]=\E[q_i]\E[k_i]=0,\qquad
  \E[Z_i^2]=\E[q_i^2]\E[k_i^2]=1,
\]
nên $\Var(Z_i)=1$. Với $i\ne j$, hai biến $Z_i$ và $Z_j$ là hàm của hai nhóm biến
độc lập $(q_i,k_i)$ và $(q_j,k_j)$, do đó $\Cov(Z_i,Z_j)=0$. Vì
$q^\top k=\sum_{i=1}^{d_k} Z_i$,
\begin{align}
  \E[q^\top k]&=0,\notag\\
  \Var(q^\top k)
    &=\sum_{i=1}^{d_k}\Var(Z_i)+2\!\!\sum_{1\le i<j\le d_k}\!\!\Cov(Z_i,Z_j)=d_k,
    \label{eq:b1-var}\\
  \Var\!\left(\frac{q^\top k}{\sqrt{d_k}}\right)
    &=\frac{1}{d_k}\Var(q^\top k)=1.\notag
\end{align}

Chú ý rằng giả thiết độc lập được dùng \emph{hai lần} với hai vai trò khác nhau:
lần thứ nhất để tách $\E[q_ik_i]$ thành tích hai kỳ vọng, lần thứ hai để triệt tiêu
hiệp phương sai giữa các chỉ số khác nhau. Chỉ giả thiết ``không tương quan'' là
không đủ cho lần thứ nhất.

\subsubsection{Vì sao điều này giúp giảm bão hòa softmax?}

Xét $m\ge2$ điểm chú ý $z_1,\ldots,z_m$ trong một hàng, với
$s_i=e^{z_i}/\sum_{j=1}^m e^{z_j}$. Dùng quy tắc đạo hàm thương, với
$\delta_{ij}=1$ nếu $i=j$ và $0$ nếu $i\ne j$:
\begin{equation}
  \frac{\partial s_i}{\partial z_j}
  =\frac{\delta_{ij}e^{z_i}}{\sum_a e^{z_a}}
   -\frac{e^{z_i}e^{z_j}}{(\sum_a e^{z_a})^2}
  =s_i(\delta_{ij}-s_j).
  \label{eq:b1-jac}
\end{equation}
Nếu điểm lớn nhất $z_a$ hơn mọi điểm khác ít nhất $g>0$ thì
\begin{equation}
  s_a=\frac{1}{1+\sum_{j\ne a} e^{-(z_a-z_j)}}
    \ \ge\ \frac{1}{1+(m-1)e^{-g}}.
  \label{eq:b1-sat}
\end{equation}
Giữ $m$ cố định và cho $g$ tăng, $s_a\to1$, các $s_j$ còn lại $\to0$, và theo
\eqref{eq:b1-jac} toàn bộ Jacobian của softmax tiến đến $0$. Khi đó trọng số chú ý
gần như không phản ứng với nhiễu nhỏ ở điểm số, tức gradient qua lớp attention bị
triệt tiêu.

Điểm then chốt: bão hòa phụ thuộc vào \emph{chênh lệch} điểm, không phụ thuộc việc
cộng một hằng số chung vào cả hàng --- softmax bất biến với phép dịch đó. Với các
tích vô hướng thô, độ lệch chuẩn của từng điểm là $\sqrt{d_k}$ theo
\eqref{eq:b1-var}. Chia cả hàng cho $\sqrt{d_k}$ đưa độ lệch chuẩn mỗi điểm về $1$
và thu nhỏ mọi chênh lệch theo cùng hệ số khi $d_k>1$.

\textbf{Phạm vi kết luận.} Kết quả phương sai giải thích \emph{lý do chọn} thang
điểm. Nó không khẳng định mọi hàng đều không bão hòa, cũng không khẳng định các tọa
độ của query và key trong mô hình đã huấn luyện là độc lập hay có phương sai đơn
vị. Giả thiết ở đây là một mô hình giả định tại khởi tạo.

%--------------------------------------------------------------------
\subsection{B.2 --- Tính đẳng biến hoán vị của self-attention}
\label{sec:b2}

\textbf{Mệnh đề.} Xét self-attention không mã hóa vị trí, không mặt nạ, dùng cùng
trọng số chiếu cho mọi hàng. Với $X\in\R^{n\times \dmodel}$ mà mỗi hàng là biểu
diễn một token, nếu $P\in\R^{n\times n}$ là ma trận hoán vị thì
$\Attn(PX)=P\Attn(X)$.

Đây là tính \emph{đẳng biến} (equivariance): đầu ra đổi thứ tự theo đầu vào. Nó
khác tính \emph{bất biến} (invariance), trong đó đầu ra không đổi.

\subsubsection{Theo dõi phép hoán vị đến ma trận điểm}

Dùng $W_Q,W_K\in\R^{\dmodel\times d_k}$ và $W_V\in\R^{\dmodel\times d_v}$, với
$n,\dmodel,d_k,d_v>0$. Đặt
\[
  Q=XW_Q,\quad K=XW_K,\quad V=XW_V,\qquad
  S=\frac{QK^\top}{\sqrt{d_k}}\in\R^{n\times n}.
\]
Theo \eqref{eq:sdpa}, $\Attn(X)=\softmax_{\mathrm{row}}(S)V$. Với đầu vào $PX$, vì
các phép chiếu tác động \emph{bên phải} còn hoán vị tác động \emph{bên trái}, ta có
$Q'=PQ$, $K'=PK$, $V'=PV$, và do đó
\[
  S'=\frac{(PQ)(PK)^\top}{\sqrt{d_k}}=P\,\frac{QK^\top}{\sqrt{d_k}}\,P^\top=PSP^\top.
\]

\subsubsection{Hoán vị đi qua softmax theo hàng}

Chọn quy ước $(PX)_{i,:}=X_{\tau(i),:}$ với $\tau$ là song ánh trên
$\{1,\ldots,n\}$. Khi đó $(PSP^\top)_{ij}=S_{\tau(i),\tau(j)}$ và
\begin{align}
  [\softmax_{\mathrm{row}}(PSP^\top)]_{ij}
    &=\frac{\exp(S_{\tau(i),\tau(j)})}{\sum_{a=1}^{n}\exp(S_{\tau(i),\tau(a)})}
     \notag\\
    &=\frac{\exp(S_{\tau(i),\tau(j)})}{\sum_{b=1}^{n}\exp(S_{\tau(i),b})}
     \label{eq:b2-key}\\
    &=[P\softmax_{\mathrm{row}}(S)P^\top]_{ij}.\notag
\end{align}
Đẳng thức \eqref{eq:b2-key} là bước quyết định: nó đúng vì $\tau$ là song ánh, nên
$a\mapsto\tau(a)$ đi qua đúng một lần mỗi chỉ số cột và tổng ở mẫu số không đổi
giá trị. Nếu $\tau$ không phải song ánh, bước này sai.

Do $P^\top P=I_n$,
\begin{align*}
  \Attn(PX)
    &=\softmax_{\mathrm{row}}(PSP^\top)\,PV
     =P\softmax_{\mathrm{row}}(S)P^\top PV\\
    &=P\softmax_{\mathrm{row}}(S)V=P\Attn(X).
\end{align*}

Trong chú ý nhiều đầu, lập luận trên đúng cho từng đầu. Ghép các đầu theo chiều đặc
trưng rồi nhân bên phải với $W_O$ vẫn cho đúng phép hoán vị hàng, vì ghép các
$PH_i$ bằng $P$ nhân với ma trận ghép các $H_i$.

\subsubsection{Giới hạn: mặt nạ nhân quả phá đẳng biến}

Với mặt nạ cộng $M\in(\R\cup\{-\infty\})^{n\times n}$ có ít nhất một vị trí không
bị che ở mỗi hàng, đặt $\Attn_M(X)=\softmax_{\mathrm{row}}(S+M)V$. Lập luận trên
vẫn cho
\begin{equation}
  \Attn_{PMP^\top}(PX)=P\Attn_M(X),
  \label{eq:b2-mask}
\end{equation}
nghĩa là \emph{nếu hoán vị cả mặt nạ} thì đẳng biến được giữ.

Điều này \emph{không} cho phép giữ nguyên một mặt nạ nhân quả cố định. Phản ví dụ:
lấy hai token vô hướng, $X=(0,2)^\top$, $W_Q=W_K=0$, $W_V=1$, và $P$ đổi chỗ hai
hàng. Vì $W_Q=W_K=0$, mọi điểm số bằng $0$, nên softmax chia đều trên các khóa
được phép. Giữ mặt nạ nhân quả (hàng một nhìn chính nó; hàng hai nhìn cả hai):
\[
  \Attn_M(X)=(0,1)^\top,\qquad
  \Attn_M(PX)=(2,1)^\top\ \ne\ (1,0)^\top=P\Attn_M(X).
\]
Vậy mệnh đề phổ quát ``attention có mặt nạ nhân quả cố định là đẳng biến'' bị bác
bỏ. Nguyên nhân cấu trúc: mặt nạ nhân quả gắn với \emph{vị trí}, không gắn với
token, nên đổi thứ tự token mà giữ mặt nạ là đổi bài toán.

Tương tự, khi cộng bảng mã hóa vị trí cố định $E$, đổi thứ tự token nhưng giữ
nguyên vị trí cho $PX+E$, nói chung khác $P(X+E)$. Vì thế định lý về attention
không PE không suy ra đẳng biến cho mô hình thực tế có PE và có mặt nạ.

\textbf{Về phản ví dụ suy biến.} Phản ví dụ dùng $d_{\mathrm{model}}=1$ và phép
chiếu bằng $0$. Điều này hợp lệ: để bác bỏ một mệnh đề phổ quát chỉ cần một ví dụ,
và kết quả dương \eqref{eq:b2-mask} vẫn được giữ nguyên bên cạnh.

%--------------------------------------------------------------------
\subsection{B.3 --- Ma trận quay khối của mã hóa vị trí hình sin}
\label{sec:b3}

\textbf{Mệnh đề.} Xét vector cột $\PE(p)\in\R^{2m}$ gồm $m\ge1$ cặp đầy đủ theo
thứ tự sin trước, cos sau: với tần số cố định $\omega_i$, $i=0,\ldots,m-1$, đặt
$\PE(p)_{2i}=\sin(p\omega_i)$ và $\PE(p)_{2i+1}=\cos(p\omega_i)$. Khi đó tồn tại
$M_\Delta\in\R^{2m\times 2m}$ chỉ phụ thuộc độ dịch $\Delta$ (với hệ tần số đã cố
định) sao cho $\PE(p+\Delta)=M_\Delta\PE(p)$ với mọi $p$. Trong Transformer gốc,
$\omega_i=10000^{-2i/\dmodel}$ \cite[§3.5]{attention}.

\subsubsection{Chứng minh cho một cặp tọa độ}

Đặt $a=p\omega_i$ và $b=\Delta\omega_i$. Công thức cộng góc cho
\begin{align*}
  \sin(a+b)&=\cos b\,\sin a+\sin b\,\cos a,\\
  \cos(a+b)&=-\sin b\,\sin a+\cos b\,\cos a.
\end{align*}
Vậy nếu $u_i(p)=(\sin(p\omega_i),\cos(p\omega_i))^\top$ thì
$u_i(p+\Delta)=M_{\Delta,i}\,u_i(p)$ với
\begin{equation}
  M_{\Delta,i}=
    \begin{pmatrix}
      \cos(\Delta\omega_i)&\sin(\Delta\omega_i)\\
      -\sin(\Delta\omega_i)&\cos(\Delta\omega_i)
    \end{pmatrix}.
  \label{eq:b3-block}
\end{equation}
Hai cột của khối này có chuẩn $1$ vì $\cos^2 b+\sin^2 b=1$; tích vô hướng của chúng
bằng $\cos b\sin b-\sin b\cos b=0$; định thức bằng $\cos^2 b+\sin^2 b=1$. Do đó
\eqref{eq:b3-block} là ma trận quay --- theo quy ước vector cột $(\sin,\cos)^\top$,
nó là phép quay góc $-b$ trong mặt phẳng tọa độ chuẩn.

\subsubsection{Ghép các cặp thành vector đầy đủ}

Lấy $M_\Delta=\diag(M_{\Delta,0},\ldots,M_{\Delta,m-1})$. Vì các khối ngoài đường
chéo bằng $0$, phép nhân tác động riêng lên từng cặp tọa độ, đúng bằng
$u_i(p)\mapsto u_i(p+\Delta)$. Ghép lại theo thứ tự ban đầu được
$M_\Delta\PE(p)=\PE(p+\Delta)$. Các phần tử của $M_\Delta$ chỉ chứa $\Delta$ và
các tần số đã cố định, \emph{không} chứa $p$ --- đây là điều cần chứng minh.

Kết quả cho thấy độ dịch cố định có \emph{biểu diễn tuyến tính} trên PE. Nó không
khẳng định mô hình đã học được cách khai thác biểu diễn đó.

\subsubsection{Phạm vi: mệnh đề cần số chiều chẵn}
\label{sec:b3-odd}

Giả thiết ``$m$ cặp đầy đủ'' là điều kiện thực sự cần, không phải quy ước trình
bày. Cài đặt \texttt{SinusoidalPositionEncoding} của nhóm chấp nhận $\dmodel$ lẻ:
với $\dmodel=2m+1$, mã sinh $m+1$ tần số, đặt $\sin$ vào $m+1$ cột chỉ số chẵn và
$\cos$ vào $m$ cột chỉ số lẻ, nên
\[
  \PE(p)=\bigl(\sin p\omega_0,\ \cos p\omega_0,\ \ldots,\
                \sin p\omega_{m-1},\ \cos p\omega_{m-1},\
                \sin p\omega_m\bigr)^\top .
\]
Tọa độ cuối không có $\cos(p\omega_m)$ đi kèm. Ta chứng minh khi đó \emph{không tồn
tại} $M_\Delta$ nào thỏa mệnh đề, trừ khi $\sin(\Delta\omega_m)=0$.

\textbf{Bổ đề.} Cho các tần số phân biệt $\omega_0,\ldots,\omega_m\in(0,\pi)$. Nếu
$L\ge 2m+2$ thì $2m+2$ hàm $p\mapsto\sin(p\omega_i)$ và $p\mapsto\cos(p\omega_i)$,
$0\le i\le m$, độc lập tuyến tính trên $\R$ khi xét như hàm trên
$\{0,1,\ldots,L-1\}$.

\emph{Chứng minh.} Đặt $z_i=e^{\mathrm i\omega_i}$. Mỗi tổ hợp tuyến tính thực của
các hàm trên viết được thành
$\sum_{i=0}^{m}\bigl(\gamma_i z_i^{\,p}+\overline{\gamma_i}\,\overline{z_i}^{\,p}\bigr)$.
Vì các $\omega_i\in(0,\pi)$ phân biệt, $2m+2$ số
$z_0,\overline{z_0},\ldots,z_m,\overline{z_m}$ đôi một khác nhau. Nếu tổng triệt
tiêu tại $p=0,\ldots,2m+1$ thì hệ Vandermonde với các nút phân biệt đó khả nghịch,
buộc mọi $\gamma_i=0$. \hfill$\square$

Với lựa chọn của Transformer gốc, $\omega_i=10000^{-2i/\dmodel}\in(0,1]\subset(0,\pi)$
và đôi một phân biệt, nên bổ đề áp dụng khi $\texttt{max\_len}\ge\dmodel+1$.

\textbf{Suy ra.} Giả sử có $M_\Delta\in\R^{(2m+1)\times(2m+1)}$ với
$\PE(p+\Delta)=M_\Delta\PE(p)$ cho mọi $p$ trong miền vị trí. Đọc hàng cuối (chỉ số
$2m$) và dùng công thức cộng góc:
\[
  \cos(\Delta\omega_m)\sin(p\omega_m)+\sin(\Delta\omega_m)\cos(p\omega_m)
  =\sum_{i=0}^{m-1}\bigl(a_i\sin(p\omega_i)+b_i\cos(p\omega_i)\bigr)
   +a_m\sin(p\omega_m),
\]
với $a_i,b_i$ là các phần tử của hàng đó, không phụ thuộc $p$. Chuyển vế:
\[
  \sin(\Delta\omega_m)\cos(p\omega_m)
  =\sum_{i=0}^{m-1}\bigl(a_i\sin(p\omega_i)+b_i\cos(p\omega_i)\bigr)
   +\bigl(a_m-\cos(\Delta\omega_m)\bigr)\sin(p\omega_m).
\]
Theo bổ đề, $\cos(p\omega_m)$ độc lập tuyến tính với mọi hàm ở vế phải, nên đẳng
thức chỉ đúng với mọi $p$ khi $\sin(\Delta\omega_m)=0$, tức
$\Delta\omega_m\in\pi\mathbb Z$. Vì $0<\omega_m\le1$, điều này loại mọi độ dịch
nguyên $\Delta$ với $|\Delta|<\pi/\omega_m$, trong đó có $\Delta=1$. Mệnh đề thất
bại ở số chiều lẻ.

\textbf{Hệ quả cho cài đặt.} Thí nghiệm A.2 dùng $\dmodel=64$ (chẵn) nên bảng PE
được dùng thực tế vẫn thỏa mệnh đề; phần kiểm thử $\dmodel$ lẻ của nhóm chỉ đối
chiếu giá trị bảng, không kiểm tính chất dịch vị trí. Muốn giữ tính chất cho số
chiều lẻ, cần bỏ tọa độ $\sin$ lẻ cặp hoặc đệm lên số chiều chẵn; đó là lựa chọn
thiết kế, không phải hệ quả của chứng minh.

%--------------------------------------------------------------------
\subsection{B.4 --- Khởi tạo, hạng, số tham số và hợp nhất LoRA}
\label{sec:b4}

\textbf{Thiết lập.} Lớp tuyến tính có $W\in\R^{d\times k}$ với $d,k\ge1$, đầu vào
cột $h\in\R^k$. LoRA đóng băng $W$, học $A\in\R^{r\times k}$ và $B\in\R^{d\times r}$
với $r\ge1$ nguyên và $\alpha$ cố định. Không huấn luyện bias; nếu có bias cố định
$b\in\R^d$, nó được cộng như nhau trong cả hai cách tính. Đặt
$\Delta W=(\alpha/r)BA$ và $W'=W+\Delta W$ \cite[§4.1--4.2]{lora}.

\subsubsection{Khởi tạo không thay đổi lớp gốc}

Khi $B=0_{d\times r}$, mỗi phần tử $(BA)_{ij}=\sum_{a=1}^{r}B_{ia}A_{aj}$ bằng $0$
bất kể giá trị của $A$. Do đó $\Delta W=0_{d\times k}$ và $W'=W$. Bài báo LoRA chọn
$A$ ngẫu nhiên và $B=0$; kết luận $W'=W$ chỉ cần điều kiện thứ hai.

Ngược lại, nếu khởi tạo \emph{cả hai} bằng $0$ thì
$\partial\mathcal L/\partial A\propto B^\top(\cdot)=0$ và
$\partial\mathcal L/\partial B\propto(\cdot)(Ah)^\top=0$, nên adapter không bao giờ
rời khỏi điểm $0$. Vì vậy một trong hai ma trận phải khác $0$ lúc khởi tạo.

\subsubsection{Hạng của cập nhật bị chặn bởi $r$}

Xem $A$ là ánh xạ $\R^k\to\R^r$ và $B$ là ánh xạ $\R^r\to\R^d$. Với mọi $h$,
$BAh=B(Ah)\in\im(B)$. Vì $\im(B)$ sinh bởi $r$ cột của $B$,
\[
  \rank(BA)=\dim\im(BA)\le\dim\im(B)=\rank(B)\le r,
\]
suy ra $\rank(\Delta W)\le r$. Mệnh đề chặn hạng của \emph{phần cập nhật}, không
chặn hạng của $W'$: tổng của một ma trận hạng đầy và một ma trận hạng thấp nói
chung vẫn hạng đầy.

\subsubsection{Đếm tham số thực sự được huấn luyện}

Học trực tiếp toàn bộ $W$ cần cập nhật $dk$ hệ số. Với LoRA, $A$ có $rk$ phần tử và
$B$ có $dr$ phần tử; vì $W$, $\alpha$ và bias không được huấn luyện, số tham số cập
nhật là
\begin{equation}
  rk+dr=r(d+k).
  \label{eq:b4-count}
\end{equation}
Số này nhỏ hơn $dk$ khi $r<dk/(d+k)$; tỷ lệ chính xác so với học cả $W$ là
$r(d+k)/(dk)$. Đây không phải tổng dung lượng mô hình, vì vẫn cần lưu trọng số gốc.

\subsubsection{Hợp nhất không thêm phép tính trong lan truyền xuôi}

Trong suy luận xác định với adapter đã cố định, hai nhánh cho cùng ánh xạ như một
lớp tuyến tính duy nhất:
\begin{equation}
  Wh+\frac{\alpha}{r}B(Ah)=\left(W+\frac{\alpha}{r}BA\right)h=W'h.
  \label{eq:b4-merge}
\end{equation}
Đây là tính phân phối của phép nhân ma trận, không phải xấp xỉ. Tính và lưu $W'$
một lần trước khi phục vụ; sau đó mỗi đầu vào dùng đúng một phép nhân
ma trận--vector $d\times k$ như lớp gốc ($dk$ phép nhân và $d(k-1)$ phép cộng nếu
tính trực tiếp, cộng cùng chi phí bias nếu có). Không còn các phép nhân phụ $Ah$,
$B(Ah)$ hay phép cộng hai nhánh trên mỗi đầu vào.

\textbf{Phạm vi kết luận.} Kết luận nói về chi phí số học \emph{sau hợp nhất}, với
giả thiết cùng kiểu thực thi lớp tuyến tính. Nó không bao gồm chi phí tạo $W'$,
không bảo đảm khớp từng bit trong số học dấu phẩy động, và không tự mở rộng thành
kết luận về độ trễ khi đổi kernel hay tái lượng tử hóa.

\subsubsection{Đối chiếu với cài đặt A.3}

Bốn tiểu mệnh đề trên khớp từng dòng với \texttt{src/lora.py}: $B$ khởi tạo bằng
\texttt{new\_zeros}; hệ số \texttt{scale = alpha / rank}; lớp nền gọi
\texttt{requires\_grad\_(False)} nên cả weight và bias bị đóng băng, đúng giả thiết
đếm tham số \eqref{eq:b4-count}; và \texttt{merge()} dùng
\texttt{addmm\_(B, A, alpha=scale)} trên bản sao, cho đúng
$W+(\alpha/r)BA$ với bias gốc, như \eqref{eq:b4-merge}. Cài đặt ràng buộc thêm
$r\le\min(d,k)$, điều mà chứng minh không cần.

%--------------------------------------------------------------------
\subsection{B.5 --- Nghiệm tối ưu có phạt KL và hàm mất mát DPO}
\label{sec:b5}

\textbf{Thiết lập.} Cố định câu hỏi $x$. Ký hiệu $R(x,y)$ là reward của đáp án $y$
(đề dùng $r(x,y)$), $\beta>0$, và $\pref(\cdot\mid x)$ là phân phối tham chiếu cố
định. Xét tập đáp án hữu hạn, không rỗng $\mathcal Y_x$, với reward hữu hạn trên
hỗ trợ $S_x=\{y\in\mathcal Y_x:\pref(y\mid x)>0\}$. Tối ưu trên mọi phân phối
$\pi(\cdot\mid x)$ có hỗ trợ nằm trong $S_x$; chính sách gán khối lượng dương cho
$y\notin S_x$ có KL vô hạn nên không thể là nghiệm hữu hạn. Quy ước: hạng tử có hệ
số xác suất bằng $0$ trong KL bằng $0$. Theo
\cite[§3--4, phương trình (1)--(7), Phụ lục A.1--A.2]{dpo}.

\subsubsection{Xây dựng một phân phối ứng viên}

Viết tắt $\pi_y=\pi(y\mid x)$, $q_y=\pref(y\mid x)$, $R_y=R(x,y)$ trên $S_x$. Hàm
mục tiêu là
\begin{equation}
  J_x(\pi)=\sum_{y\in S_x}\pi_y R_y
    -\beta\!\!\sum_{y\in S_x:\,\pi_y>0}\!\!\pi_y\log\frac{\pi_y}{q_y}.
  \label{eq:b5-obj}
\end{equation}
Đặt $Z(x)=\sum_{y\in S_x}q_y\exp(R_y/\beta)$. Vì tập hữu hạn, hỗ trợ không rỗng và
các số hạng dương hữu hạn, $0<Z(x)<\infty$. Định nghĩa
\begin{equation}
  \pi^*(y\mid x)=\frac{\pref(y\mid x)\exp(R(x,y)/\beta)}{Z(x)}
  \quad(y\in S_x),
  \label{eq:b5-gibbs}
\end{equation}
và $\pi^*(y\mid x)=0$ ngoài $S_x$. Tổng tử số trên $S_x$ chính bằng $Z(x)$, nên
$\pi^*$ là một phân phối, dương trên toàn bộ $S_x$.

\subsubsection{Chứng minh tối ưu toàn cục và tính duy nhất}

Ký hiệu $\pi_y^*=\pi^*(y\mid x)$. Từ \eqref{eq:b5-gibbs},
$\log(\pi_y^*/q_y)=R_y/\beta-\log Z(x)$. Thế vào và dùng $\sum_y\pi_y=1$:
\begin{align}
  \beta\KL(\pi\|\pi^*)
     &=\beta\!\!\sum_{y:\,\pi_y>0}\!\!\pi_y\log\frac{\pi_y}{\pi_y^*}\notag\\
     &=\beta\!\!\sum_{y:\,\pi_y>0}\!\!\pi_y\log\frac{\pi_y}{q_y}
       -\sum_y\pi_y R_y+\beta\log Z(x)\notag\\
     &=\beta\log Z(x)-J_x(\pi).
  \label{eq:b5-central}
\end{align}
Đẳng thức \eqref{eq:b5-central} là trục của chứng minh: cực đại hóa $J_x$ tương
đương cực tiểu hóa $\KL(\pi\|\pi^*)$, vì $\beta\log Z(x)$ không phụ thuộc $\pi$.

Để thấy điều kiện dấu bằng mà không cần KKT, dùng $-\log t\ge 1-t$ với $t>0$, dấu
bằng chỉ tại $t=1$. Với $T=\{y:\pi_y>0\}$, áp dụng cho $t=\pi_y^*/\pi_y$:
\begin{equation}
  \KL(\pi\|\pi^*)\ \ge\ \sum_{y\in T}(\pi_y-\pi_y^*)
       \ =\ 1-\sum_{y\in T}\pi_y^*\ \ge\ 0.
  \label{eq:b5-unique}
\end{equation}
Nếu $\KL=0$, bất đẳng thức đầu buộc $\pi_y^*=\pi_y$ trên $T$, và bất đẳng thức sau
buộc $\pi^*$ không gán khối lượng ngoài $T$. Ngược lại, $\pi=\pi^*$ cho $\KL=0$.
Vì $\beta>0$, hàm $J_x$ đạt cực đại $\beta\log Z(x)$ \emph{đúng tại} $\pi^*$; không
còn nghiệm nào khác, kể cả ở biên của miền phân phối.

\textbf{Miền đếm được.} Một điều kiện đủ để dùng nguyên lập luận là reward bị chặn
hai phía trên hỗ trợ: nếu $a\le R_y\le b$ thì
$0<e^{a/\beta}\le Z(x)\le e^{b/\beta}<\infty$, nên kỳ vọng reward hữu hạn và
$\log(\pi_y^*/q_y)$ bị chặn. Đẳng thức \eqref{eq:b5-central} vẫn đúng khi KL nhận
giá trị $+\infty$, cho mục tiêu $-\infty$ chứ không tạo hiệu $\infty-\infty$. Nếu
reward không bị chặn, phải kiểm riêng việc chuẩn hóa và tính xác định của mục tiêu.

\subsubsection{Từ Bradley--Terry tới xác suất ưu tiên}

Xét cặp $y_w,y_l\in S_x$ cho cùng $x$, trong đó $y_w$ được ưu tiên. Mô hình
Bradley--Terry gán
\begin{align}
  \Pr(y_w\succ y_l\mid x,y_w,y_l)
   &=\frac{\exp(R(x,y_w))}{\exp(R(x,y_w))+\exp(R(x,y_l))}\notag\\
   &=\sigma\bigl(R(x,y_w)-R(x,y_l)\bigr),
   \qquad \sigma(z)=\frac{1}{1+e^{-z}}.
  \label{eq:b5-bt}
\end{align}
Giải \eqref{eq:b5-gibbs} theo reward:
\begin{equation}
  R(x,y)=\beta\log\frac{\pi^*(y\mid x)}{\pref(y\mid x)}+\beta\log Z(x).
  \label{eq:b5-invert}
\end{equation}
Lấy hiệu giữa hai đáp án \emph{cùng} $x$, hai hạng $\beta\log Z(x)$ triệt tiêu. Do
đó xác suất ưu tiên chỉ còn phụ thuộc hiệu hai log-tỷ số chính sách tối ưu trên
tham chiếu. Bước triệt tiêu \emph{không} áp dụng cho hai đáp án của hai câu hỏi
khác nhau --- đây là lý do dữ liệu ưu tiên phải ghép theo cùng một prompt.

\subsubsection{Hợp lý cực đại theo chính sách tham số hóa}

Dùng $\pi_\theta$ thay cho $\pi^*$ chưa biết trong \eqref{eq:b5-bt}. Hai đáp án của
mỗi cặp phải có xác suất dương dưới \emph{cả} $\pi_\theta$ và $\pref$ để log-tỷ số
hữu hạn. Đặt
\begin{equation}
  z_\theta(x,y_w,y_l)
   =\beta\left[
      \log\frac{\pi_\theta(y_w\mid x)}{\pref(y_w\mid x)}
      -\log\frac{\pi_\theta(y_l\mid x)}{\pref(y_l\mid x)}
   \right].
  \label{eq:b5-logit}
\end{equation}
Với tập $\mathcal D$ các bộ ba đã sắp theo nhãn ưu tiên, mỗi quan sát đóng góp xác
suất $\sigma(z_\theta)$ vào hợp lý của nhãn, có điều kiện trên câu hỏi và cặp đáp
án. Với các quan sát độc lập có điều kiện, hợp lý là tích các xác suất đó; lấy log,
đổi dấu và lấy trung bình trên $\mathcal D$:
\begin{equation}
  \boxed{\ \mathcal L_{\mathrm{DPO}}(\theta)
    =-\E_{(x,y_w,y_l)\sim\mathcal D}\bigl[\log\sigma(z_\theta(x,y_w,y_l))\bigr].\ }
  \label{eq:b5-loss}
\end{equation}

Không cần tính $Z(x)$ và không cần huấn luyện reward model riêng để tính mất mát.
Kiểm dấu: nếu hai log-tỷ số trong \eqref{eq:b5-logit} bằng nhau thì mất mát một cặp
bằng $-\log(1/2)=\log 2\approx0{,}693$; nếu hiệu tăng thì $\sigma$ tăng và mất mát
giảm.

\textbf{Phân biệt hai loại duy nhất.} Tính duy nhất ở \eqref{eq:b5-unique} là duy
nhất của $\pi^*$ \emph{với reward đã cho} khi tối ưu trên các phân phối. Nó không
phải tính duy nhất của tham số $\theta$ khi học từ hữu hạn cặp.

\subsubsection{Đối chiếu với cài đặt A.4}

Biểu thức \eqref{eq:b5-logit} khớp đúng \texttt{src/dpo.py}: mã tính
\texttt{margin = (pw - pl) - (rw - rl)} rồi \texttt{-logsigmoid(beta * margin)
.mean()}, tức $-\log\sigma(z_\theta)$ lấy trung bình trên các cặp, đúng
\eqref{eq:b5-loss}. Ràng buộc $\beta$ hữu hạn và dương trong mã tương ứng với giả
thiết $\beta>0$ ở đây.

%--------------------------------------------------------------------
\subsection{B.6 --- Cross-entropy, KL và perplexity}
\label{sec:b6}

\textbf{Thiết lập.} Phân phối dữ liệu cố định $\pdata$ và họ phân phối mô hình
$p_\theta$ trên cùng miền rời rạc hữu hạn $\mathcal Y$. Giả sử $p_\theta(y)>0$ với
mọi $y$ sao cho $\pdata(y)>0$ khi cần giá trị hữu hạn. Hạng tử có $\pdata(y)=0$
được bỏ khỏi tổng; nếu $\pdata(y)>0$ nhưng $p_\theta(y)=0$ thì cả cross-entropy và
KL đều là $+\infty$.

\subsubsection{Phân rã cross-entropy kỳ vọng}

Đặt $S=\{y:\pdata(y)>0\}$. Theo định nghĩa,
\begin{align*}
  H(\pdata,p_\theta)&=-\sum_{y\in S}\pdata(y)\log p_\theta(y),\\
  H(\pdata)&=-\sum_{y\in S}\pdata(y)\log\pdata(y).
\end{align*}
Khai triển tỷ số trong KL thành hiệu hai logarit:
\begin{equation}
  \KL(\pdata\|p_\theta)
   =\sum_{y\in S}\pdata(y)\log\frac{\pdata(y)}{p_\theta(y)}
   =-H(\pdata)+H(\pdata,p_\theta).
  \label{eq:b6-decomp}
\end{equation}
Vì $H(\pdata)$ hữu hạn và không phụ thuộc $\theta$, hai mục tiêu khác nhau đúng một
hằng số: chúng xếp thứ tự các mô hình như nhau và có cùng tập nghiệm cực tiểu (nếu
cực tiểu đạt được):
\begin{equation}
  \argmin_\theta H(\pdata,p_\theta)=\argmin_\theta \KL(\pdata\|p_\theta).
  \label{eq:b6-equiv}
\end{equation}
Trên miền đếm được, lập luận trừ hằng số vẫn dùng được khi $H(\pdata)$ hữu hạn.
Không được lấy hiệu $\infty-\infty$ nếu entropy dữ liệu vô hạn.

\subsubsection{Phiên bản có ngữ cảnh cho dự đoán token}

Ký hiệu $C$ là ngữ cảnh và $Y$ là token đích. Giữ phân phối của $C$ cố định, áp
dụng \eqref{eq:b6-decomp} cho mỗi ngữ cảnh $c$ rồi lấy kỳ vọng:
\begin{align}
  \E_{(C,Y)\sim\pdata}[-\log p_\theta(Y\mid C)]
   &= H_{\pdata}(Y\mid C)\notag\\
   &\quad+\E_{C\sim\pdata}\bigl[\KL(\pdata(\cdot\mid C)\|p_\theta(\cdot\mid C))\bigr],
  \label{eq:b6-ctx}
\end{align}
trong đó $H_{\pdata}(Y\mid C)=\E_{C}[H(\pdata(\cdot\mid C))]$ không phụ thuộc mô
hình và hữu hạn vì từ vựng hữu hạn. Đây là phát biểu tương ứng khi mô hình ngôn ngữ
được huấn luyện dự đoán token tiếp theo theo ngữ cảnh đúng từ dữ liệu.

\subsubsection{Từ log-xác suất trung bình theo token đến PPL}

Trên tập đánh giá, gọi $(c_t,y_t)$ là ngữ cảnh và token đích ở vị trí được chấm thứ
$t$, tổng cộng $N>0$ token. Với logarit tự nhiên và xác suất dương tại các vị trí
này,
\begin{equation}
  \bar H=-\frac{1}{N}\sum_{t=1}^{N}\log p_\theta(y_t\mid c_t),\qquad
  \PPL=e^{\bar H}
    =\left(\prod_{t=1}^{N} p_\theta(y_t\mid c_t)\right)^{-1/N}.
  \label{eq:b6-ppl}
\end{equation}
Đẳng thức cuối dùng $\exp(\sum_t\log a_t)=\prod_t a_t$. Vậy PPL là nghịch đảo
trung bình nhân của xác suất token đúng; vì hàm mũ tăng chặt, cross-entropy theo
token thấp hơn tương đương PPL thấp hơn \emph{trên cùng thiết lập chấm}.

\emph{Ví dụ minh họa.} Hai token có xác suất đúng $1/2$ và $1/8$ cho
$\bar H=\tfrac12(\log 2+\log 8)=\log 4$ và $\PPL=4$.

\subsubsection{Quy tắc tổng hợp và phân biệt với DPO}
\label{sec:b6-agg}

Khi cộng qua nhiều batch hoặc nhiều câu: cộng tổng log-xác suất âm rồi chia tổng số
token được chấm; \emph{không} lấy trung bình các PPL riêng lẻ. Lý do: $\PPL$ là hàm
mũ lồi của $\bar H$, nên trung bình các PPL không bằng PPL của trung bình (bất đẳng
thức Jensen), và các batch nói chung có số token khác nhau. Loại các vị trí padding
khỏi cả tử và mẫu. Khi so sánh, dùng cùng tokenizer, dữ liệu và quy ước chấm token.

Ngược lại, $\log\pi_\theta(y\mid x)$ trong \eqref{eq:b5-logit} là log-xác suất
\emph{cả câu trả lời}. Theo quy tắc dây chuyền nó bằng
$\sum_{t=1}^{T}\log\pi_\theta(y_t\mid x,y_{<t})$ với $T$ là số token của đáp án
theo cùng quy ước kết thúc chuỗi cho hai mô hình. Không tính log-xác suất của prompt
cố định, và \emph{không} chia tổng này cho $T$ trong DPO chuẩn.

Lý do của sự khác biệt: PPL đo \emph{chất lượng trên mỗi token} nên phải chuẩn hóa
theo độ dài để so sánh được giữa các tập; còn $\log\pi_\theta(y\mid x)$ trong DPO
là xác suất của \emph{một lựa chọn} trong mô hình Bradley--Terry
\eqref{eq:b5-bt}, và chia cho $T$ sẽ thay đổi bài toán so sánh ưu tiên theo độ dài
đáp án.

\subsubsection{Đối chiếu với cài đặt A.2}

Quy tắc ở Mục~\ref{sec:b6-agg} khớp \texttt{validation\_metrics} trong
\texttt{src/char\_lm.py}: mã tích lũy \texttt{nll\_sum} ở \texttt{float64} với
\texttt{reduction='sum'}, chia cho \texttt{token\_count} rồi mới lấy
\texttt{math.exp}. Docstring ghi rõ ``không lấy trung bình PPL các batch''.
